\documentclass[11pt,letterpaper]{article}

\usepackage[margin=1in]{geometry}
\usepackage{enumitem}
\usepackage{hyperref}
\usepackage{amsmath}
\usepackage{amssymb}

\title{Spacecraft ADCS Homework 3}
\author{Due Thursday, April 2, 2026}
\date{}

\begin{document}

\maketitle

\section{Attitude Sensors Revisited}

We will now refine and expand on the sensor models you implemented in the last homework. We will categorize the sensors into three types, and implement realistic error models for each:
\begin{enumerate}
	\item \textbf{Vector or bearing sensors:} This includes things like magnetometers, sun sensors, and horizon sensors that we looked at previously. For sensors in this category, you should implement an affine error model with additive noise. The model will have the form:
		\begin{equation*}
			y = M y_{\mathrm{ideal}} + b + w
		\end{equation*}
		where $M$ is a constant matrix containing scale and misalignment errors, $b$ is a constant bias vector, and $w$ is additive noise sampled from a Gaussian. Choose reasonable values for $M$, $b$ and $W = \mathrm{cov}(w)$ that you justify based on evidence from the literature.
	\item \textbf{Star trackers:} Star trackers generally return a full attitude quaternion as a measurement. They typically have much larger errors about the boresight vector compared to the other two directions. These errors should be modeled as random rotations applied to each quaternion measurement and can be sampled from a Gaussian defined over some three-parameter attitude parameterization (e.g. axis-angle vectors).
	\item \textbf{Gyroscopes:} Gyros utilize a similar affine error model to vector sensors, except that the bias is time-varying and is typically modeled by a random walk:
		\begin{equation*}
			y_\mathrm{gyro} = M \omega_{\mathrm{true}} + b(t) + w
		\end{equation*}
		Select reasonable values for $M$ and the covariances associated with both $w$ and the random-walk process for $b(t)$ that will be modeled in your MEKF. Justify these choices with references to the literature.		
\end{enumerate}

\section{Recursive Attitude Estimation}

Implement an MEKF like the one we discussed in class. Test it using simulated time history data generated using your spacecraft dynamics model and the sensor models you have developed. For now, simulations should be performed in a slow (few degrees per second) uncontrolled tumble. Test using many random initial conditions. You should investigate and document the following:
\begin{enumerate}
	\item What sample rate should you run your filter at? This should be based on your mission requirements (maximum slew or spin rate).
	\item How do the errors from your filter compare to the static estimates computed previously?
	\item Is your estimator consistent? Make sure that the covariance estimates generated by the filter are consistent with the actual state errors.
	\item Explore the convergence behavior of your filter by initializing it with different state and covariance values. How good must your initial guess be? How long does it take the filter to converge to its steady state accuracy?
\end{enumerate}


\end{document}
